% Verrerie usuelle du laboratoire de chimie (annexe « Fiches expérimentales »)
\begin{tikzpicture}[x=1cm,y=1cm,scale=0.7,every node/.style={transform shape=false}]
\pic at (0,0) {becher={largeur=1.9,hauteur=2.3,niveau=1.2,liquide=pcEau}};
\pic at (3.0,0) {erlen={niveau=0.7,liquide=pcEau}};
\pic at (6.0,0) {fiole={hauteur=3.4,niveau=1,liquide=pcEau}};
\pic at (9.1,0) {pipette={hauteur=4.6,niveau=1,liquide=pcEau}};
\pic at (12.4,0.9) {burette={hauteur=5.4,niveau=4.4,liquide=pcEau}};
\foreach \x/\t in {0/{bécher},3.0/{erlenmeyer},6.0/{fiole jaugée},9.1/{pipette jaugée},12.4/{burette graduée}}{
  \node[pcLegendeG,anchor=north] at (\x,-0.25) {\t};}
\end{tikzpicture}
\par\smallskip\legende{Le bécher et l'erlenmeyer contiennent et mélangent ; la fiole et la pipette jaugées mesurent un volume précis ; la burette délivre un volume variable, lu sur ses graduations.}
